% ============================================================
%  EASRP 2026, just-module-creator
% ============================================================

\documentclass[11pt,a4paper]{article}

\usepackage{graphicx}
\usepackage{easrp2026}              % anonymous submission
% \usepackage[accepted]{easrp2026}  % camera-ready version
\usepackage{amsmath}
\usepackage{amssymb}
\usepackage{booktabs}
\usepackage{hyperref}
\usepackage{url}

% Author details are used only in the accepted version.
\easrpauthor{Anonymous}
\easrpaffiliation{Withheld for review}

\easrptitle{Creating and Sharing Genomic Annotation Modules with AI: The
Just-DNA Ecosystem}

\begin{document}

\maketitle

\begin{abstract}
Genome annotation depends on curated knowledge that changes as papers and
databases publish new variant-trait associations. The Just-DNA ecosystem stores
this knowledge in genomic annotation modules: versioned data packages that
describe which genotype a claim concerns, what evidence supports it, and how a
consumer should present it. A module can be written and edited manually. An AI
assistant can also draft one from a paper, a variant list, or a description of
a topic, but its output remains an editable data artifact rather than generated
analysis code.

We present the Just-DNA ecosystem and the path a module follows through it.
\texttt{just-dna-format} defines the schema; \texttt{just-dna-enricher} records
source-derived facts; \texttt{just-dna-compiler} validates and compiles the
files. \textbf{just-module-creator}, an open-source plugin for Claude Code and
Codex, is the authoring entry point: it exposes these packages as typed tools so
that both AI-assisted and manually edited modules use the same live schema. An
author can keep the result in a local module store, rehearse publication in a
deletable test registry, or publish a version to the production registry.
Just-DNA-Lite later filters a local VCF and joins selected modules to matching
genotypes. Polygenic scores follow a parallel numerical path through
\texttt{just-prs} (described in a companion platform paper).

The language model helps with reading and drafting, while deterministic software
handles schema validation, compilation, genome processing, and scoring. This
paper describes the module lifecycle from source material to shared store and
local use. It also defines a protocol for evaluating variant recovery, citation
identity, and effect-direction agreement in AI-assisted creation. The software
is open-source and intended for research use only.
\end{abstract}

% ================================================================
\section{Introduction}
% ================================================================

Suppose a researcher wants a recent lactose-tolerance study to inform a genomic
report. The researcher can transcribe the relevant variants and evidence into
tables, or ask an AI assistant to prepare a first draft. In either case, the
result is the same kind of object: a module that can be edited, checked,
versioned, and shared. The researcher may keep it in a local store while it is
being reviewed, publish it to a test registry for rehearsal, and later release
an accepted version through the production registry. When a user selects that
module in Just-DNA-Lite, the application joins its records to matching
genotypes in a local VCF and includes the resulting annotations in a report.

This separation between knowledge and sample processing is central to the
design. A paper may report an association with lactose tolerance, drug response,
or another trait, but a computational pipeline needs an unambiguous variant,
the genotype to which the claim applies, the genome build and allele
orientation, a structured conclusion, and a trace back to the evidence. A
\emph{module} is a versioned collection of those records. It contains no sample
data and no executable analysis code. The same compiled module can be applied
to many genomes.

The sample follows a different path. A VCF often contains millions of variant
calls. Just-DNA-Lite normalizes their representation, applies configured
quality filters, and joins the retained calls to selected modules and reference
data. Polygenic risk scores follow a parallel numerical path: \texttt{just-prs}
applies a PGS Catalog model to the sample and places the score in the context of
a reference population. These maintained pipelines can
reuse normalized columnar data instead of reparsing the original VCF for every
question.

A general-purpose assistant is useful during curation, but it does not provide
this structure by itself. Asked to ``make me a lactose-tolerance module,'' it
may return rsIDs, an explanation, and a new script that searches a VCF. The
answer can contain a fabricated identifier, a reversed allele, or a real PMID
that names the wrong paper \citep{ji2023survey,gao2023enabling}. Generated code
may also bypass the established filters and validators or depend on an unstated
environment. The next run can produce a different implementation because the
implementation was generated as part of the answer.

The Just-DNA ecosystem gives that assistant a narrower and more useful job.
\texttt{just-module-creator} helps it read sources, prepare module files, run
checks, and present unresolved decisions to the author. The files use schemas
from \texttt{just-dna-format}; \texttt{just-dna-enricher} records facts from
external sources; \texttt{just-dna-compiler} validates and compiles them; and
\texttt{just-dna-registry} stores published versions. Manual authors use the
same files and tools. Just-DNA-Lite and \texttt{just-prs} remain responsible for
sample-level computation.

This paper follows a module through the ecosystem, from manual or AI-assisted
authoring to a shared store and local use in a genomic report. Just-DNA-Lite
and \texttt{just-prs} are described as the consumer side of the boundary; their
internal design is the subject of a companion platform paper.

This paper makes four contributions:

\begin{enumerate}
    \item It describes the module contract shared by manual editors, AI-assisted
    authoring tools, versioned stores, and the genome-analysis consumer.
    \item It presents \texttt{just-module-creator}, an authoring plugin that
    exposes the public capabilities of \texttt{just-dna-format},
    \texttt{just-dna-compiler}, \texttt{just-dna-enricher}, and
    \texttt{just-dna-registry} as typed MCP tools. The plugin reads the
    installed schema instead of carrying a copy.
    \item It defines how modules move between a local store, a deletable test
    registry, and the production registry while preserving provenance and
    decisions that still require review.
    \item It defines an evaluation protocol for variant recovery, citation
    identity, and effect-direction agreement. Compiler round-trip behaviour is
    tested by \texttt{just-dna-compiler} and is not counted again as module
    creation accuracy.
\end{enumerate}

% ================================================================
\section{Related work}
% ================================================================

\paragraph{Genome annotation and catalogs.}
A variant call records what differs between a sample and a reference genome.
Annotation adds information about that call, such as its predicted consequence,
known clinical classification, reported trait association, or the conclusion
assigned to a particular genotype. VEP, ANNOVAR, and OpenCRAVAT/OakVar perform
this kind of lookup by joining a callset to established resources
\citep{mclaren2016vep,wang2010annovar,pagel2020opencravat}. ClinVar and the GWAS
Catalog publish records that can supply annotation content
\citep{landrum2018clinvar,sollis2023gwas}. The problem addressed here begins
before the join: selected evidence must be turned into a reviewable,
machine-checkable collection that can be distributed and applied consistently.

\paragraph{Language models and scientific extraction.}
Language models can produce fluent text around a false identifier or citation
\citep{ji2023survey}. Tool use reduces the need to improvise code when a model
can call a typed operation instead \citep{schick2023toolformer,anthropic2024mcp}.
Biomedical MCP services can search literature, ontologies, and variant
databases. just-module-creator uses the same general approach, then carries the
result into module files that the rest of the Just-DNA pipeline can inspect and
reuse.

\paragraph{Agent orchestration.}
Many scientific-agent systems assign research and review to several model
instances. The plugin described here does not require its own orchestration
runtime. Claude Code or Codex supplies the agent, while the plugin supplies the
tools and the written workflow. This paper evaluates that authoring path, not a
particular choice of language model or agent-team topology.

% ================================================================
\section{Method}
% ================================================================

\subsection{How modules move through the ecosystem}

Module authoring and sample analysis remain separate until Just-DNA-Lite joins
a compiled module to a local genome. On the authoring path, a person may edit
the files directly or use an AI assistant to draft them. Both routes then pass
through the same enrichment and compilation steps before the module enters a
local or shared store.

\begin{table}[ht]
\centering
\small
\caption{The module lifecycle and the two sample-processing paths.}
\label{tab:paths}
\begin{tabular}{@{}p{0.16\linewidth}p{0.77\linewidth}@{}}
\toprule
Path & Stages \\
\midrule
Module lifecycle & papers and databases $\rightarrow$ manual editing or AI
draft $\rightarrow$ enrichment and compilation $\rightarrow$ local store or
registry \\
Sample annotation & VCF $\rightarrow$ normalization and quality filtering
$\rightarrow$ joins with selected modules and reference data $\rightarrow$
annotated tables and report \\
Polygenic scoring & VCF plus a PGS Catalog model $\rightarrow$
\texttt{just-prs} computation $\rightarrow$ score and reference-population
context $\rightarrow$ report \\
\bottomrule
\end{tabular}
\end{table}

On the sample path, normalization gives equivalent variants a consistent
representation and computes the genotype used by later joins. Configured
quality filters can remove calls that fail VCF filter status, depth, or quality
criteria. The pipeline writes the normalized data to Parquet, so each selected
module can use a streaming join instead of independently parsing the original
VCF. An optional Ensembl join adds reference annotations. Report generation
reads the resulting module-specific Parquet files and their display metadata.

Polygenic risk scores are part of the ecosystem but not calculations performed
by \texttt{just-module-creator}. Just-DNA-Lite delegates them to
\texttt{just-prs}, which applies published PGS Catalog scoring models to the
sample. A module can describe a published score in \texttt{pgs.csv}; the
numerical score still belongs to \texttt{just-prs}. This keeps literature
curation and sample-level computation separate while allowing both results to
appear in the same application.

Within \texttt{just-module-creator}, the language model participates only in the
knowledge path. It can search and read sources, draft module rows, and help
review disagreements. An assistant may launch Just-DNA-Lite or
\texttt{just-prs} through their established interfaces, but it does not replace
their implementations. VCF filtering, module joins, and polygenic scoring
remain ordinary software steps with explicit inputs and repeatable code.

\subsection{The module artifact}

A just-dna module begins as a directory. It always contains
\texttt{module\_spec.yaml} and at least one recognized table. The table records
the thing being annotated, such as a genotype, a diplotype, or a measured
quantity. Evidence and machine-produced sidecars are stored separately. The
compiler turns this directory into parquet files and a content-addressed
\texttt{manifest.json} \citep{justdnaformat2026}.

The authored files are ordinary YAML, CSV, Markdown, and optional image files.
They can be created and revised without an AI assistant. just-module-creator
operates on this same directory rather than introducing an AI-specific module
format. The module describes how a genotype or measurement should be
interpreted if a consumer encounters it. Just-DNA-Lite, or another compatible
consumer, supplies the genome or measurement later.

The table kind follows the subject of the claim. For example, in format 0.6.6
used for this manuscript, a single-variant claim belongs in
\texttt{variants.csv}; evidence for those rows belongs in
\texttt{studies.csv}. Pharmacogenomic diplotypes, copy-number ranges, repeat
counts, and published polygenic-score identifiers use other tables. The
manuscript does not reproduce their column lists. The plugin calls
\texttt{describe\_table} and \texttt{table\_requirements}, which read the
installed pydantic models. An author therefore sees the schema accepted by the
compiler running in that session.

\subsection{Components and responsibilities}

No single package implements the whole path. Each component owns a boundary
that another component can call without reimplementing its rules. Within the
module toolchain, dependencies point inward: enricher $\rightarrow$ compiler
$\rightarrow$ format. The format and compiler do not fetch remote data. The
enricher performs source lookups and writes the derived sidecars used later by
the compiler. The registry publishes the compiled result, and Just-DNA-Lite
consumes it in a separate application. Polygenic scoring is delegated to
\texttt{just-prs}.

\begin{table}[ht]
\centering
\caption{Responsibilities across the Just-DNA ecosystem.}
\label{tab:packages}
\begin{tabular}{@{}p{0.25\linewidth}p{0.36\linewidth}p{0.30\linewidth}@{}}
\toprule
Component & Responsibility & Does not \\
\midrule
\texttt{just-module-creator} & research and authoring workflow & process a genome \\
\texttt{just-dna-format} & schema and identity rules & access the network \\
\texttt{just-dna-compiler} & validate, compile, reverse, and sign & access the network \\
\texttt{just-dna-enricher} & resolve, draft, and cross-check & decide scientific meaning \\
\texttt{just-dna-registry} & publish, discover, and download modules & author module rows \\
Just-DNA-Lite & filter VCFs, join annotations, and produce reports & define module schema or compiler rules \\
\texttt{just-prs} & compute polygenic scores & create annotation modules \\
\bottomrule
\end{tabular}
\end{table}

Of these components, just-module-creator is the agent-facing authoring entry
point examined most closely in this paper. It is distributed as an application
whose public contract is the MCP tool surface and the accompanying workflow
documents. The same source tree ships plugin manifests for Claude Code and
Codex. Both hosts launch the same server and load the same skills.

\subsection{Tools available to the assistant}

The MCP tools are grouped by task: authoring, research, checks, enrichment,
comparison, and registry work. The default configuration lists the full
surface. An optional layered configuration initially lists the core authoring
loop and a \texttt{toolbox} command that reveals the other groups. Tool search
can reduce the listing further. These discovery options change what the agent
sees in its context; they do not create a separate implementation of the
workflow.

Registry writes require a token, with one necessary exception:
\texttt{registry\_register} creates the token and therefore cannot require one
in advance. Tools whose cost depends on a large source corpus state that cost
in their descriptions.

Literature and enrichment requests share a pacing gate, including the NCBI
budget used by several lookups. The literature search fills a gap left by the
enricher. The enricher can verify a PMID that an author already has, but it does
not search for the paper. \texttt{literature\_search} returns titles so the
assistant can check identity. The existence of a PMID alone says nothing about
whether it is the intended article.

\subsection{How a person enters the workflow}

A person usually knows what they want to accomplish, not which internal stage
performs it. They may want to create a module, understand an inherited
directory, find a paper, revise a published module, or decode an error message.
The command menu uses those intentions as its entries.

\begin{center}
\small
\begin{tabular}{@{}ll@{}}
\toprule
Command & Purpose \\
\midrule
\texttt{/create-module} & begin or continue creating a module \\
\texttt{/module-status} & inspect a directory and identify the next decision \\
\texttt{/module-revise} & begin another pass over an existing module \\
\texttt{/find-evidence} & find and read evidence for a row \\
\texttt{/module-publish} & rehearse and then publish \\
\texttt{/module-symptom} & explain a message from the toolchain \\
\texttt{/module-install-local} & install a compiled module without a catalog \\
\bottomrule
\end{tabular}
\end{center}

The first version of the menu exposed twenty stage names. That forced a person
to choose between terms such as \texttt{module-curate} and
\texttt{module-enrich} before the system had explained either one. The seven
commands now route to thirteen guides. A guide contains the procedure for one
stage and is loaded by path when the assistant reaches that stage.

\texttt{/create-module} is therefore a router. It examines what the author has
already provided and selects the next stage. If the request shows that the
person first needs an explanation, the router loads the overview guide. The
overview is not a menu command because a new author would not know its internal
name.

\subsection{The authoring journey}

The workflow follows the module from an idea to one of its destinations:

\begin{center}
\small
\texttt{origin $\rightarrow$ scaffold $\rightarrow$ draft $\rightarrow$ curate
$\rightarrow$ enrich $\rightarrow$ check $\rightarrow$ compile $\rightarrow$
close $\rightarrow$ rehearse $\rightarrow$ publish}
\end{center}

The origin determines where the work begins. A source such as ClinVar, CPIC, or
ClinPGx may already publish rows, in which case a drafting tool can prepare a
partial table. A paper or a free-text request instead begins with evidence
search and manual authoring. Scaffolding creates the files and placeholders
required for the selected table kind.

Curation follows drafting because a drafted row may still contain decisions
that the source cannot make for the module. The author or assistant must decide
which genotype the claim concerns, how a weight should be interpreted, and how
the conclusion should be worded. The workflow presents these decisions rather
than silently choosing values that merely satisfy the schema.

Enrichment resolves identifiers and writes sidecars such as
\texttt{resolution.csv} and \texttt{literature.csv}. These files record the
source answers used for later offline compilation. Compilation itself does not
use the network. The plugin always calls the compiler with
\texttt{resolve\_with\_ensembl=True}. The parameter name is misleading: setting
it to false disables all resolution, including an injected
\texttt{resolution.csv}, and the compiler can then succeed with null
coordinates that cannot match a VCF.

After compilation, \texttt{close} records a statement in
\texttt{verification.json} and binds it to the authored files. Editing those
files invalidates the closure. In the current 0.x format a module can still
compile and publish without a closure, but it carries a warning. Closing is
an attributed declaration of completion. It does not prove that the module is
correct.

Most modules return to this workflow. A later source release, a new paper, or a
reviewer's correction begins a revision pass. The router inspects the existing
directory so that the assistant does not treat a revision as a blank first
draft and overwrite decisions whose history matters.

\subsection{Edits, evidence, and responsibility}

The compiler reports problems but does not edit authored values.
just-module-creator is the authoring application, so its workflow permits
edits. That permission comes with limits.

First, tool-mediated overrides append a record to
\texttt{logs/authoring.log} and preserve a reason in
\texttt{provenance.json}. A general hand edit is not captured automatically.
The server therefore instructs an assistant to call
\texttt{record\_override} after such a change. This is a known gap between the
desired provenance record and what the current tools can enforce.

Second, the workflow does not fill a value from the source that will later be
used to check that same value. Such a check would compare the source with
itself. The same problem appears when a row uses an article title as its
\texttt{provenance\_quote}: the title is guaranteed to occur in the article, so
the resulting match provides no evidence that anyone located support for the
row's claim.

Third, disagreement with an archive requires review of both sides. An archive
may lag a retraction, a meta-analysis, or a later reclassification. Replacing
the module's value with the archive's value can therefore make a module worse.
\texttt{record\_override} stores which field differs, what the source said, who
made the decision, and why the authored value should remain. The record does
not turn the disagreement into a passed check.

An assistant may read retrieved full text and locate a
\texttt{provenance\_quote}. It must copy the passage verbatim and identify who
located it in \texttt{StudyRow.curator}. This attribution helps a reviewer
decide where to look closely. It does not transfer responsibility away from
the human author.

\subsection{Local and shared module stores}

In this paper, a module store is a place from which a compiled module can be
installed or discovered. A local installation is useful during authoring. The
registry provides two shared stores for publication rehearsal and release.

The registry has two independent instances. Production is the catalog used by
consumers. The polygon, selected with \texttt{target=test}, is a rehearsal
environment where an author can delete a test publication. Writes default to
the polygon. Catalog reads require an explicit target so that an author does
not publish to one instance, read from the other, and mistake the missing
result for a failed publication.

The production catalog is immutable. Publishing a version also claims its
authored content by hash, and yanking the version does not release that claim.
The workflow therefore rehearses on the polygon and asks separately before a
production publish. The registry accepts the specification, runs its own
enrichment and strict compilation, and stores the artifact it produced. It
does not rely on a locally claimed digest, and it does not require the author
to sign the local artifact.

Publication is not the only way to use a module.
\texttt{/module-install-local} installs a compiled artifact into a
Just-DNA-Lite checkout so it can be tested with a local VCF. This route tests
the annotation connection between the module and the consumer. A polygon
publication tests the separate connection between the module and the registry.

% ================================================================
\section{Results}
% ================================================================

\subsection{Implemented surface}

just-module-creator is released in the 0.24 series at the time of this draft.
The same source tree loads in Claude Code and Codex. After \texttt{uv sync},
the tools named by the authoring workflow are registered. The former
essentials and extended modes are gone; the optional layered listing only
changes how tools are revealed to the agent.

Schema answers in the environment used for this manuscript come from
\texttt{just-dna-format} and \texttt{just-dna-compiler} 0.6.6. The plugin does
not maintain a separate list of columns or vocabularies.

The command redesign reduced the text placed in every session from 14,688
characters for twenty entries to 3,464 characters for seven entries, a 76\%
reduction. A test walks the links from those seven commands and confirms that
every guide remains reachable. The smaller menu therefore removes repeated
prompt content without making a stage inaccessible.

\subsection{Compiler behaviour belongs to the compiler}

Within one compiler version, \texttt{just-dna-compiler} tests that a valid
specification can be compiled, reversed, and compiled again without changing
its authored content signature. just-module-creator calls that compiler and
fixes the resolution setting described above. It does not implement a second
compiler or present the upstream round-trip tests as evidence that an assistant
extracted a paper correctly.

This separation avoids a circular evaluation. If an assistant writes a wrong
claim in valid syntax, the compiler may preserve it perfectly. Reproducible
output shows that the compiler kept the input stable; it does not show that the
input agrees with the literature.

\subsection{Protocol for creation accuracy}

The creation evaluation begins with an expert-checked fixture. Each fixture
contains a free-text prompt and the \texttt{module\_spec.yaml},
\texttt{variants.csv}, and \texttt{studies.csv} that the assistant should
recover. The generated module is compared with the fixture on three questions:

\begin{itemize}
    \item Variant recall is the fraction of ground-truth
    $(\textit{rsID},\textit{genotype})$ rows present in the generated module.
    \item Citation identity requires every cited PMID to exist and to name the
    paper indicated by the title returned during search. Existence alone is
    insufficient.
    \item Weight-sign accuracy measures whether the generated effect direction
    agrees with the fixture. Magnitude is reported separately because a module
    weight is an authored choice and is not copied from a GWAS beta.
\end{itemize}

The development fixtures \texttt{fto\_bmi}, which contains one locus at
rs1421085, and \texttt{longevity\_2026} are too small to support a performance
claim. A larger adjudicated set and repeated runs are still needed. We
therefore report no recall or precision estimate in this draft.

\subsection{How to read a green result}

Several outputs look reassuring while answering a narrower question than a
reader may expect. Strict compilation checks reproducibility and applies the
compiler's blocking rules; it does not establish biological correctness. A
digest match likewise shows that bytes or authored content agree under a
specified comparison.

A source timeout produces an unknown result, not a negative one and not a
pass. The tools preserve this distinction with null and unknown values. The
quote example is another useful test of a green check. If every row cites its
article title as the supporting passage, the full-text matcher will find every
quote even though nobody located evidence for any individual claim.

The creation protocol counts none of these as evidence of accurate extraction.
It remains a protocol until it has been run on a sufficiently large fixture set
using this plugin.

% ================================================================
\section{Discussion}
% ================================================================

The ecosystem's architecture changes what an AI assistant is asked to produce.
It writes module data for an existing pipeline instead of inventing the
pipeline again for each question. The same compiler and schema are used for
every module, and the sample is normalized once before selected annotations are
joined through the established columnar implementation. The language model may
still misunderstand a paper. Its draft is kept inspectable, and source
disagreements and unresolved judgements remain visible to the author.

The workflow must work for people with different levels of expertise. A
geneticist may want to decide every genotype and evidence statement. A
non-specialist may provide a topic and several papers, then rely on the
assistant to explain each decision. The tool surface cannot assume which kind
of author is present. It can perform mechanical corrections and keep a record
of them, while presenting scientific interpretation and production publishing
for review.

The skills help the agent follow this path, but they are instructions rather
than a security boundary. An agent can still bypass them and edit a CSV
directly. The same is true of the provenance log: tool-mediated overrides are
recorded, while a hand edit depends on the agent calling the recording tool.
These limits are important because the paper evaluates a workflow composed of
tools and instructions, not a closed application that prevents every other
action.

Agent hosts and MCP conventions change quickly. This paper focuses on the parts
that should survive a client update. Schema answers come from the installed
models. Edits can carry attribution, and source disagreements remain visible.
Publication is rehearsed on a separate instance. The compiler retains authority
over the artifact. Client-specific setup belongs in the repository
documentation.

The present evaluation has several limitations. It contains no creation
accuracy estimate, and the available fixtures cover too few kinds of module.
The plugin cannot determine whether an annotation is medically correct. It can
make the claim structured, attributable, and easier to review. The module
authoring workflow targets GRCh38 and does not perform liftover. The wider
ecosystem does compute polygenic scores through \texttt{just-prs}, but that
numeric path is not evaluated in this paper. Local installation into
Just-DNA-Lite is described as a procedure that runs in that checkout;
just-module-creator does not execute the consumer on the author's behalf.
We also have not compared the runtime of freely generated VCF scripts with the
Just-DNA-Lite pipeline. The performance argument made here is architectural:
the established path reuses normalized Parquet data and a maintained join
implementation instead of generating a new analysis for each request.

% ================================================================
\section{Conclusion}
% ================================================================

The Just-DNA ecosystem separates two kinds of work. Module authors turn papers
and databases into structured, attributable annotation knowledge. Just-DNA-Lite
filters and normalizes a local genome, joins it to selected compiled modules,
computes polygenic scores through \texttt{just-prs}, and produces results for
inspection and reporting. A companion platform paper describes that consumer
side in detail.

Within the authoring side, just-module-creator lets an AI assistant draft and
revise ordinary module files while the installed schema, enricher, compiler,
and registry provide the shared path from draft to published artifact. A module
can then be installed locally or consumed from the catalog without asking the
language model to write another genomic pipeline. Scientific judgement belongs
to the author, and sample-level computation belongs to deterministic software.
The software is open-source and intended for research use only.

% References
\bibliographystyle{plainnat}
\bibliography{references}

% Appendix
\appendix
\section{Code and data availability}
\label{app:code}

\begin{itemize}
    \item \textbf{just-module-creator}:
    \url{https://github.com/dna-seq/just-module-creator}
    (MCP server, Claude Code and Codex plugin manifests, skills; MIT).
    \item \textbf{just-dna-format, just-dna-compiler, and just-dna-enricher}:
    \url{https://github.com/dna-seq/just-dna-format}
    (schema, reference compiler, and enricher).
    \item \textbf{just-dna-registry}: production catalog at
    \url{https://module-registry.just-dna.life} and polygon at
    \url{https://module-polygon.just-dna.life}.
    \item \textbf{Just-DNA-Lite}:
    \url{https://github.com/dna-seq/just-dna-lite}
    (local VCF processing, annotation, and reporting; described in a
    companion platform paper).
    \item \textbf{just-prs}:
    \url{https://github.com/dna-seq/just-prs}
    (polygenic risk-score computation library).
\end{itemize}

\section{Research use only}
\label{app:ruo}

Annotation modules summarize published association findings. They are not
clinical-grade evidence. The authoring plugin does not make individual-level
predictions and does not open a genome. Just-DNA-Lite and \texttt{just-prs}
process sample data for research and educational use; their reports and scores
are not diagnoses or medical advice. Language that implies causation, such as
``causes'' or ``guarantees,'' is outside the scope of a module written with this
plugin.

\end{document}
